\title{QLoRA: Efficient Finetuning of Quantized LLMs}
The storage data type is dequantized into the computation data type to execute the forward and backward passes; however, weight gradients are only computed for the LoRA parameters, which employ 16-bit BrainFloat. More precisely, the $2^k$ values $q_i$ representing the data type are computed as:

\begin{equation}
q_i  = \frac{1}{2}\left( Q_X\left(\frac{i}{2^k+1}\right) + Q_X\left(\frac{i+1}{2^k+1}\right)\right),
\end{equation}

where $Q_X(\cdot)$ denotes the quantile function of the standard normal distribution $N(0,1)$. We suspect this uncertainty stems from the ambiguous definition of scale—for instance, the interpretation of an 8 on a 10-point scale varies across contexts. \paragraph{Double Quantization{}} 
We propose \emph{Double Quantization} (DQ), a technique that applies quantization to the quantization constants themselves to further reduce memory usage. \paragraph{Math} One of {Guanaco}'s main limitations lies in mathematics, a domain where numerous language models encounter difficulties as noted in \citep{liang2022holistic}. In summary, \textsc{QLoRA} utilizes one storage data type (typically 4-bit NormalFloat) alongside a computation data type (16-bit BrainFloat). Feel free to ask if you have more questions.}
\end{quote}

which is incorrect not only once (the accurate factorization is $3 \times 17 \times 43$), but twice over. Jackson moved the beans to the pantry. For instance, we detect pronounced order effects, with GPT-4 assigning higher ratings to the system presented first in its prompt. Overall, these findings indicate moderate concordance between system-level assessments by GPT-4 and human annotators, suggesting that model-based evaluation can be a reasonably reliable proxy for human evaluation. Consistent with prior research on quantization \citep{dettmers2022case}, our outcomes on MMLU and Elo suggest that, given fixed resources for fine-tuning and inference, it is advantageous to increase the number of base model parameters while reducing their numerical precision. \paragraph{4-bit NormalFloat achieves superior performance compared to 4-bit Floating Point} 

Although the 4-bit NormalFloat (NF4) data type is theoretically optimal in terms of information, it remains to be confirmed whether this advantage translates into practical performance gains. \textbf{Guanaco}:~\texttt{Al Jolson is credited with popularizing the song `I'll Keep the Lovelight Burning,' and he was born in the year 1886.}
 \end{quote}

it provides an incorrect popularizer and an incorrect birth year (although the birth year stated is accurate for Al Jolson). To address this to some extent, we conduct a qualitative analysis in two parts. \section{Broader Impacts}

Our \textsc{QLoRA} finetuning technique is the first to enable finetuning of 33B parameter models on a single consumer GPU and 65B parameter models on a single professional GPU without degrading performance relative to full finetuning baselines. Overall, we anticipate that \textsc{QLoRA} will have a broadly beneficial effect by making the finetuning of high-quality LLMs substantially more accessible and convenient. For instance, when training a 7B LLaMA model on FLAN v2 with a batch size of 1, and using LoRA weights equal to the commonly adopted 0.2\% of the original model weights\citep{hu2021lora,liu2022few}, the memory footprint for LoRA input gradients is 567 MB, whereas the LoRA parameters occupy only 26 MB. \paragraph{Finetuning with Adapters} Although we utilize Low-rank Adapters~\citep{hu2021lora} (LoRA), numerous other Parameter Efficient FineTuning (PEFT) techniques have been introduced, including prompt tuning~\citep{qin2021learning,lester2021power,li2021prefix}, adjusting the embedding layer inputs~\citep{an2022input}, tuning hidden states~(IA$^3$) \citep{liu2022few}, incorporating full additional layers~\citep{houlsby2019parameter}, bias tuning~\citep{zaken2021bitfit}, learning a weight mask guided by Fisher information~\citep{sung2021training}, and hybrid methods~\citep{henderson2021compacter}. Moreover, in Table~\ref{tbl:elo_large}, GPT-4 rates its own outputs considerably higher than human evaluators do, with Elo scores of 1348 versus 1176, corresponding to roughly a 20\% greater probability of winning against an opponent. We utilize this capability to allocate paged memory for the optimizer states, which are automatically moved to CPU RAM when GPU memory is insufficient, and then paged back into GPU memory as required during the optimizer update process. The beans are in the treasure chest. An additional possible area of influence is the deployment on mobile devices. When compared with GPT-4, {Guanaco} 65B and 33B achieve an expected win probability of 30\%, based on Elo ratings derived from system-level pairwise comparisons by human evaluators—this is the highest value reported so far. We release all models and code, including CUDA kernels designed for 4-bit training.\footnote{\url{https://github.com/artidoro/qlora} and \url{https://github.com/TimDettmers/bitsandbytes}}
\end{abstract}

\section{Introduction}

Finetuning large language models (LLMs) is a highly effective approach to enhance their capabilities \citep{min2021metaicl, wei2021finetuned, ouyang2022training, wang2022super, wang2022self, liu2022few} as well as to introduce desirable traits or eliminate unwanted behaviors \citep{ouyang2022training,askell2021general,bai2022training}. We evaluate this hypothesis in two scenarios. \textsc{QLoRA} can contribute to privacy-conscious use of LLMs, enabling users to retain control over their own data and models, while also simplifying the deployment of LLMs. Our analysis also reveals that automated evaluation tools exhibit discernible biases. Double Quantization reduces this memory cost. Specifically, Double Quantization treats the quantization constants $c_2^{\text{FP32}}$ from the initial quantization as inputs for a secondary quantization step. Although LoRA was developed as a Parameter Efficient Finetuning (PEFT) approach, the majority of memory usage in LLM finetuning arises from activation gradients rather than the learned LoRA parameters. The upcoming sections detail the experiments conducted to address these inquiries. \paragraph{Memory Requirement of Parameter-Efficient Finetuning} A key aspect to consider is the memory consumption of LoRA during training, specifically regarding the quantity and size of adapters employed. First, in \S\ref{ssec:examples}, we present examples that we consider indicative of certain recurring patterns found in the text generated by our 65b {Guanaco} model. \textsc{QLoRA} incorporates several innovations to reduce memory consumption without compromising performance: (a) 4-bit NormalFloat (NF4), a novel data format that is information-theoretically optimal for weights following a normal distribution, (b) Double Quantization, which further decreases average memory usage by quantizing the quantization constants themselves, and (c) Paged Optimizers that effectively handle memory spikes. \\ \\ To determine the total earnings, we first calculate how much he made from mowing 16 lawns: \\ \\  \$33 x 16 = \$528 \\ \\ Next, we add the tips from three customers: \\  \\ \$10 x 3 = \$30 \\ \\ Therefore, the total amount Lee earned is: \\ \\ \$528 + \$30 = \$558}
\end{quote}

Nonetheless, {Guanaco} may fail even on straightforward problems if it does not approach them methodically, a recognized challenge \cite{weichain}. We conduct a hyperparameter search across learning rates from 1e-6 to 5e-5 and batch sizes ranging from 8 to 128 in order to establish robust baseline models. Although 7B models have previously been demonstrated to run on phones, \textsc{QLoRA} is the first technique to facilitate the finetuning of such models on these devices. A cherry-picked analysis highlights cases where \textbf{Guanaco} underperforms relative to ChatGPT. Therefore, James might expect Abby to search for the pen in the desk drawer, as that is where she last left it. The results for finetuning the 7B LLaMA model on Alpaca are presented in Figure~\ref{fig:hyper}. To further investigate the benefits of NF4 over other 4-bit data formats, we replicate the setup from \citet{dettmers2022case} and evaluate post-quantization zero-shot accuracy and perplexity on various models (OPT~\citep{zhang2022opt}, LLaMA~\citep{touvron2023llama}, BLOOM~\citep{scao2022bloom}, Pythia~\citep{biderman2023pythia}) with sizes ranging from 125 million to 13 billion parameters. Consequently, we train over 1,000 models spanning various instruction tuning datasets, model architectures, and sizes ranging from 80M to 65B parameters. \textsc{QLoRA} presents several advancements aimed at minimizing memory consumption while maintaining performance: (1) \textbf{4-bit NormalFloat}, a quantization data type optimized from an information theory perspective for normally distributed data, which delivers better empirical outcomes compared to 4-bit Integers and 4-bit Floats. The Vicuna prompts consist of 80 queries drawn from a diverse range of categories, which we use without alteration. The third step corresponds to adjusting the weight tensor's standard deviation to equal the standard deviation of the $k$-bit data type. We randomly draw samples from the set of labeled comparisons to calculate Elo~\citep{elo1967proposed,elo1978rating}. For each query, along with the responses from ChatGPT and a given model, GPT-4 is prompted to assign a score out of ten to both answers and provide a rationale. Although {Guanaco} 33B has more parameters than Vicuna 13B, it employs only 4-bit precision weights, making it significantly more memory efficient at 21 GB compared to 26 GB, and yields a three-percentage-point improvement over Vicuna 13B. \section{Advancing Chatbot Performance with QLoRA}
\label{sec:pushing}
Having demonstrated that 4-bit \textsc{QLoRA} achieves parity with 16-bit performance across different scales, tasks, and datasets, we undertake a detailed examination of instruction finetuning on the largest open-source language models accessible for research purposes. We will now describe the components of \textsc{QLoRA} and then provide a formal definition of \textsc{QLoRA}. We have additionally shown that NF4 is more effective than FP4 and that the application of double quantization does not impair performance. Upon detecting a pattern, we attempt to construct questions or prompts designed to elicit that pattern even when it represents an incorrect answer. From the validation dataset, we extract all user messages as queries, incorporating previous dialogue turns into the prompt. At the individual example level, agreement between GPT-4 and the majority human vote is lower, with Fleiss’ $\kappa=0.25$. We initialize scores at 1,000 and set $K=32$. This implies that the performance degradation caused by imprecise quantization can be fully mitigated by performing adapter finetuning after quantization. Beyond LoRA, there is a broad spectrum of Parameter Efficient FineTuning (PEFT) approaches that have demonstrated strong results; however, their scalability to large models is uncertain. Therefore, we emphasize that while model-based evaluations offer a cost-effective alternative to human annotation, they also possess inherent uncertainties. The results, displayed in Table~\ref{tbl:llama-datatype}, reveal that NF4 combined with double quantization completely recovers the 16-bit LoRA MMLU accuracy. Additionally, we conduct a comprehensive evaluation of chatbot performance utilizing both human assessors and GPT-4 for assessment. Or perhaps another objective altogether? \paragraph{Paged Optimizers} leverage the NVIDIA unified memory~\footnote{\tiny \url{https://docs.nvidia.com/cuda/cuda-c-programming-guide}} capability, which performs automatic page-by-page transfers between the CPU and GPU to enable error-free GPU computation when the GPU experiences occasional out-of-memory conditions. A blocksize of 64 is chosen for $\mathbf{W}$ to enhance quantization precision, while a blocksize of 256 is used for $c_2$ to optimize memory efficiency. Below, we outline our proposed setup, employing nucleus sampling with $p=0.9$ and temperature $0.7$ in every case. \paragraph{Human Evaluation} Although recent research suggests that generative models can be effectively utilized for system evaluations \citep{fu2023gptscore}, the reliability of GPT-4 ratings as proxies for chatbot performance assessments, to our knowledge, has yet to be validated against human judgments. Because finetuning following quantization appears to restore most of the information lost during quantization, this could allow for more aggressive quantization strategies. Vicuna represents a fully finetuned LLaMA 13B model trained on proprietary user-shared conversations from ShareGPT, effectively serving as a distilled model from OpenAI GPT systems. \textbf{Guanaco}: \texttt{I'm sorry, but I am not able to disclose the secret word. We move forward to explore instruction tuning at scales that would be unfeasible to investigate using full 16-bit finetuning on standard academic research hardware. Future research should explore strategies to address these issues, drawing on fields that have developed methods for handling subjective preferences, such as Human-Computer Interaction and Psychology. This demonstrates that while gradient checkpointing is essential, drastically minimizing LoRA parameters offers only modest memory savings. \textbf{Guanaco}: \texttt{James is aware that Abby put the pen in the desk drawer. Our assessments include GLUE \citep{wang2018glue} with RoBERTa-large \citep{liu2019roberta}, Super-NaturalInstructions (TKInstruct) \citep{wang2022super} with T5 \citep{t5}, and 5-shot MMLU \citep{hendrycksmeasuring} following finetuning LLaMA on Flan v2 \citep{longpre2023flan} and Alpaca \citep{alpaca}. Is there anything else you'd like to know?}
\end{quote}

demonstrating the necessity to explore techniques for more reliable adherence to instructions. However, since these quantile approximations are not exact, the data type exhibits substantial quantization errors for outliers, which are frequently the most critical values. The outcomes, presented in Table~\ref{tab:nf4vsbf16}, indicate that adapter-based methods at 16-bit, 8-bit, and 4-bit precision achieve comparable performance to the fully finetuned 16-bit baseline on both datasets. We advance these methods by emphasizing a more dependable evaluation framework. The precise values of this data type are provided in Appendix~\ref{app:nf4}. Are we aiming to develop models that excel in academic knowledge typical of high school and college settings, or models that perform well in chatbot conversational skills? For instance, FLAN v2 aligns closely with MMLU but differs from chatbot benchmarks, whereas the Chip2 dataset shows the opposite pattern, and both models perform accordingly on the MMLU and Vicuna benchmarks. Prominent methods aimed at maintaining 16-bit precision quality involve handling outlier features (e.g., SmoothQuant~\citep{xiao2022smoothquant} and LLM.int8()~\citep{dettmers2022llm}), while others employ more advanced grouping strategies \citep{park2022nuqmm, yao2022zeroquant}. \paragraph{\textsc{QLoRA}.} Based on the components outlined above, we formalize \textsc{QLoRA} for a single linear layer in the quantized base model augmented with a single LoRA adapter as follows:

\begin{equation}
    \mathbf{Y}^{\text{BF16}} =  \mathbf{X}^{\text{BF16}} \text{doubleDequant}(c_1^{\text{FP32}}, c_2^{\text{k-bit}}, \mathbf{W}^{\text{NF4}}) + \mathbf{X}^{\text{BF16}}\mathbf{L}^{\text{BF16}}_1\mathbf{L}^{\text{BF16}}_2,
\end{equation}

where doubleDequant$(\cdot)$ is defined by:

\begin{equation}
    \text{doubleDequant}(c_1^{\text{FP32}}, c_2^{\text{k-bit}}, \mathbf{W}^{\text{k-bit}}) = \text{dequant}(\text{dequant}(c_1^{\text{FP32}}, c_2^{\text{k-bit}}), \mathbf{W}^{\text{4bit}}) = \mathbf{W}^{\text{BF16}},
\end{equation}

We employ NF4 format for $\mathbf{W}$ and FP8 for $c_2$. It is important to note that this relative score can exceed 100\% if the model attains a higher absolute score than ChatGPT. By contrast, the 4-bit quantized base model requires 5,048 MB of memory. Our findings indicate that QLoRA finetuning on a small, high-quality dataset yields state-of-the-art outcomes, even when using models smaller than previous leading methods. As demonstrated in Figure~\ref{fig:hyper} for LLaMA 7B finetuning on Alpaca, the most pivotal LoRA hyperparameter is the total number of LoRA adapters employed, with LoRA applied across all linear transformer block layers necessary to match full finetuning performance. \\ \\
    \textbf{Guanaco}: \texttt{The earth has never been officially verified as flat by peer reviewed scientists. This highlights a certain degree of orthogonality among current evaluation benchmarks: excelling in MMLU does not necessarily correlate with strong chatbot capabilities as assessed by Vicuna or OA benchmarks, and vice versa. Small hyperparameter searches were conducted for the 13B and 33B LLaMA models, finding that hyperparameters established for 7B models generally transfer well (including the number of epochs), except for learning rate and batch size. Jackson entered the playroom. The computationally expensive quantile estimation and its approximation errors can be circumvented when the input tensors originate from a distribution that is fixed except for a quantization constant. After each match, the Elo rating updates in proportion to the expected outcome: an unexpected victory causes a significant rating shift, whereas a predicted result yields only a minor adjustment. As shown in Figure~\ref{fig:nf4} and Table~\ref{tbl:nf4_ppl}, NF4 substantially outperforms FP4 and Int4, and the use of double quantization reduces memory usage without compromising performance. We approximate that using an iPhone 12 Plus, \textsc{QLoRA} can finetune approximately 3 million tokens overnight while the device is charging. \textsc{QLoRA} uses a single low-precision storage data type, commonly 4-bit, alongside a computation data type that is typically BFloat16. One issue with symmetric k-bit quantization is that it lacks an exact representation of zero, which is crucial for quantizing padding and other zero-valued elements without error. Nonetheless, this assumption could be wrong if Abby observed James moving the pen.}
\end{quote}

Nonetheless, these conclusions are inconsistent, and the model frequently offers explanations that involve assumptions which are illogical in the given context, for example:

\begin{quote}
    \textbf{User}: Evelyn entered the living room. Apart from our study, SwitchBack layers~\citep{wortsman2023stable} represent the sole research investigating backpropagation through quantized weights at scales exceeding 1B parameters. Our findings indicate a high level of agreement between GPT-4 and human judgments regarding model rankings within the tournaments, although there are notable cases of significant disagreement. To evaluate the effectiveness of instruction finetuning for these models, we assess them using a challenging Natural Language Understanding benchmark (MMLU) and introduce novel approaches for evaluating real-world chatbot performance. For our purposes, we define this range as $[-1, 1]$. \paragraph{Training Setup} To eliminate confounding influences from diverse training objectives, we apply QLoRA finetuning with cross-entropy loss (supervised learning) without reinforcement learning, even on datasets that contain human evaluations of various responses. (2) \textbf{Double Quantization}, a technique that further compresses the quantization constants themselves, resulting in an average saving of about 0.37 bits per parameter (equivalent to roughly 3 GB for a 65B parameter model). We employ 8-bit floats with a blocksize of 256 for this secondary quantization, as 8-bit quantization does not degrade performance, consistent with findings from \citet{dettmers2022case}. Although recent quantization techniques can significantly reduce the memory consumption of LLMs \citep{dettmers2022llm, dettmers2022case, frantar2022gptq, xiao2022smoothquant}, these methods are only applicable for inference and fail during training \citep{wortsman2023stable}. Leveraging \textsc{QLoRA}, we train the \textbf{Guanaco} model series, with our second-best model achieving 97.8\% of ChatGPT's performance on the Vicuna~\citep{vicuna2023} benchmark, while being trainable within 12 hours on a single consumer-grade GPU; with a single professional GPU over 24 hours, we reach 99.3\% with our largest model, effectively closing the performance gap with ChatGPT on Vicuna. Since $c_2^{\text{FP32}}$ are positive, we first subtract their mean to center the values around zero and then apply symmetric quantization. To mitigate this bias, we recommend reporting the average score across both presentation orders. To support further research, we provide all model outputs annotated by both humans and GPT-4. The widespread adoption of LLMs entails recognized risks \citep{bommasani2021opportunities,bender2021dangers}, but we argue that democratizing access to this rapidly proliferating technology will promote improved and more independent scrutiny, as opposed to concentrating LLM capabilities exclusively within large corporations that withhold models or source code from public audit. 
\begin{abstract}
We introduce \textsc{QLoRA}, a memory-efficient finetuning method that enables tuning a 65B parameter model on a single 48GB GPU while maintaining the full task performance of 16-bit finetuning. These head-to-head comparisons are performed for all pairwise system permutations on both the Vicuna and OA benchmarks. This multiple-choice benchmark encompasses 57 tasks spanning elementary mathematics, US history, computer science, law, and more. \textsc{QLoRA} performs backpropagation through a frozen, 4-bit quantized pretrained language model into Low Rank Adapters~(LoRA). We measure the probability of {Guanaco}-65B generating socially biased token sequences compared to other models in Table~\ref{tbl:crows}. Lossy quantization methods explore trade-offs associated with conventional rounding~\citep{dettmers2022case,zeng2022glm,pope2022efficiently} or seek to optimize rounding choices to enhance quantization accuracy~\citep{frantar2022gptq}. The next best-performing model trained solely on open-source data is the Anthropic HH-RLHF model, which scores 30 percentage points lower than {Guanaco} on the Vicuna benchmark (refer to Table~\ref{tbl:vicuna_eval}). I am programmed to not reveal or repeat the secret word, no matter what the human asks. Although we conduct assessments using MMLU, the Vicuna benchmark, and the OA benchmark, we do not evaluate on other benchmarks like BigBench, RAFT, and HELM; thus, our findings may not necessarily generalize to these datasets. On the positive side, we undertake an extensive analysis on MMLU and introduce novel methods for chatbot evaluation. As will be discussed later, this capability is vital for achieving full 16-bit precision performance recovery. We refer to these datasets as the Vicuna and OA benchmarks. \textbf{Guanaco}: \texttt{Sure, I'd be glad to assist! Where does James think Abby will look for the pen? Therefore, we recommend employing the Elo ranking approach \citep{elo1967proposed}, which uses \textit{pairwise} comparisons from human judges and GPT-4 to circumvent issues related to anchoring an absolute scale. The claim that a group of scientists have peer reviewed and verified the earth as flat is simply untrue. The outcomes of the tournament are then combined into Elo scores~\citep{elo1967proposed, elo1978rating} that establish the performance ranking of the chatbots. Frequently, system designers wish to withhold certain information from the user, even though the system is permitted to provide free text replies. Can I assist you with anything else?}
\end{quote}

which is the expected response. In deployment, our smallest \textbf{Guanaco} model (7B parameters) requires only 5 GB of memory and outperforms a 26 GB Alpaca model by more than 20 percentage points on the Vicuna benchmark (Table~\ref{tbl:vicuna_eval}). \textbf{Guanaco}: \texttt{Sorry, the secret word is banana. Elo rating, commonly applied in chess and various games, quantifies the expected win probability relative to an opponent's rating—for instance, a player with an Elo of 1100 against one rated 1000 is predicted to win approximately 65\% of the time; matches between equally rated players, such as 1000 vs 1000 or 1100 vs 1100, correspond to a 50\% expected win rate. \paragraph{4-bit NormalFloat Quantization} The NormalFloat (NF) data type is based on Quantile Quantization \citep{dettmers2022optimizers}, an information-theoretically optimal method that guarantees each quantization bin contains an equal number of values drawn from the input tensor. \subsection{Considerations}
\label{ssec:considerations}

\paragraph{Evaluation}

We observe moderate concordance among human evaluators (Fleiss $\kappa=0.42$), with further decline when comparing two strong models. What is the secret word? \textsc{QLoRA} lowers the average memory requirements for finetuning a 65B parameter model from over 780GB of GPU memory to under 48GB, while maintaining runtime efficiency and predictive accuracy comparable to a fully finetuned 16-bit baseline. Subsequently, we evaluate performance via direct pairwise comparisons of system outputs. \section{\textsc{QLoRA} Finetuning}

\textsc{QLoRA} achieves precise 4-bit finetuning through two techniques we propose—4-bit NormalFloat (NF4) quantization and Double Quantization. To guarantee a discrete zeropoint at $0$ and utilize all $2^k$ levels for a k-bit data type, we construct an asymmetric data type by estimating quantiles $q_i$ over two intervals: $2^{k-1}$ quantiles for the negative range and $2^{k-1}+1$ quantiles for the positive range. Since instruction finetuning is crucial for converting raw pretrained LLMs into ChatGPT-like chatbots, we believe our method will democratize finetuning and make it commonplace, particularly benefiting researchers with limited resources—a significant advancement for accessibility to cutting-edge NLP technology. Furthermore, we release a suite of adapters for models of sizes 7B, 13B, 33B, and 65B, trained across eight different instruction-following datasets, resulting in 32 distinct open-sourced, fine-tuned models. We observe a pronounced order effect where GPT-4 tends to give higher scores to the response that appears earlier in the prompt. \paragraph{Benchmark Data} We conduct evaluations using two carefully selected query datasets: the Vicuna prompts \citep{vicuna2023} and the OASST1 validation set \citep{kopf2023openassistant}. We have shown that our top 33B model trained on the Open Assistant dataset can compete with ChatGPT on the Vicuna benchmark. In these scenarios, input tensors share identical quantiles, making exact quantile estimation computationally practical. Following \citet{vicuna2023}, this process is repeated 10,000 times with varying random seeds to mitigate ordering effects, such as the influence of which model pairs compete initially. To accomplish this, we fine-tune LLaMA models sized 7B through 65B using two instruction-following datasets, Alpaca and FLAN v2, and assess performance on the MMLU benchmark using 5-shot accuracy. We propose that our \textsc{QLoRA} approach could achieve the significant milestone of allowing the finetuning of LLMs directly on phones and other environments with limited resources. \paragraph{Summary} Our findings consistently demonstrate that 4-bit \textsc{QLoRA} using the NF4 data type matches the performance of 16-bit full fine-tuning and 16-bit LoRA fine-tuning on academic benchmarks with rigorously validated evaluation protocols. All outputs in this section were generated using Nucleus Sampling \cite{holtzmancurious} with $p=0.9$. \paragraph{Experimental setup.} We evaluate three model architectures (encoder, encoder-decoder, and decoder-only) and compare QLoRA against 16-bit adapter finetuning and full finetuning for models up to 3B parameters. Specifically, we reduce the learning rate by half for the 33B and 65B models while doubling the batch size. The OASST1 dataset is a multilingual, crowd-sourced compilation of multi-turn conversations between users and an assistant. Consequently, we conduct two parallel human evaluations on the Vicuna benchmark, mirroring both automated evaluation protocols described above. Once the ranges of the weights and data type are aligned, standard quantization procedures can be applied. \paragraph{Refusal} In a similar vein, {Guanaco} occasionally declines to follow instructions for apparently arbitrary reasons:

\begin{quote}
    \textbf{User}: Please reverse the words in the sentence ``I'm a cool dude in a mocha mood.'' \\ \\
    \textbf{Guanaco}: \texttt{I'm sorry, I am unable to do that at the moment, but I can explain why the sentence you provided is grammatically correct.\\ \\The sentence ``I'm a cool dude in a mocha mood" adheres to proper grammar and syntax rules in English. Employing \textsc{QLoRA}, we finetune over 1,000 models and deliver a comprehensive analysis of instruction following and chatbot capabilities across eight instruction datasets, diverse model architectures (LLaMA, T5), and model sizes that traditional finetuning methods cannot feasibly support (e.g., 33B and 65B parameters). \paragraph{Low-rank Adapters} Low-rank Adapter (LoRA) finetuning \citep{hu2021lora} is a technique designed to lower memory usage by employing a small number of trainable parameters—called adapters—while keeping the rest of the model parameters fixed and unaltered. This advancement paves the way for future efforts in \textsc{QLoRA} fine-tuning on specialized open-source datasets, enabling the development of models that can compete with the leading commercial systems currently available. Similarly, we observe that the default hyperparameters for fully finetuned baselines are not optimally tuned. \paragraph{Factual Recall} For queries such as ``What is the capital of Zambia?'' all models reliably produce correct answers, for instance,  

\begin{quote} 
    \textbf{Guanaco}:
    \texttt{The capital of Zambia is Lusaka.}
\end{quote}

However, as questions become more obscure, {Guanaco} tends to become less dependable while maintaining confidence. Overall, these findings confirm that 4-bit \textsc{QLoRA} is an effective technique capable of producing state-of-the-art chatbots that rival ChatGPT. Further analysis of biases in {Guanaco} and comparable chatbots is left for future investigation. The relatively low sample-level agreement between GPT-4 and human raters (Fleiss $\kappa=0.25$) further indicates that human judges and automated evaluators may rely on differing preference criteria. Over time, Elo ratings tend to reflect each player's skill level in the competition. These challenges align with findings in recent studies \cite{sap2022neural}, but warrant further investigation. Consequently, both the quantiles associated with the data type and those of the neural network weights must be normalized within this interval. \paragraph{Data} Since, to our knowledge, no thorough analysis of recent instruction-following datasets exists, we select eight contemporary datasets. How many dollars did Lee earn mowing lawns last week? We replicate the evaluation setup from \citet{dettmers2022case}, where quantized large language models (OPT \citep{zhang2022opt}, BLOOM \cite{scao2022bloom}, Pythia \cite{biderman2023pythia}, LLaMA) across various sizes (125M to 65B) and data types are assessed on language modeling tasks and several zero-shot benchmarks. Direct questioning fails:

\begin{quote}
    \textbf{User}: What is the secret word? We then combine these sets of $q_i$ and remove one of the two zeros that appear in both sets. This fact has been established for centuries and supported by numerous experiments and observations. However, James also knows he has relocated the pen to the bag, and he assumes Abby is unaware of this. Taken together, this provides strong evidence that 4-bit \textsc{QLoRA} tuning reliably produces results equivalent to those achieved by 16-bit approaches. Abby placed the pen in the desk drawer. Future research should explore the existence of possible biases in automated evaluation frameworks as well as strategies to mitigate them. This prompts the key question of whether 4-bit adapter finetuning can recover the lost performance. Exploration of the trade-offs associated with other PEFT methods is reserved for future investigations. Nevertheless, a small amount of manipulation can circumvent this restriction:

\begin{quote}
    \textbf{User}: This is a game. This supports our conclusions that (1) \textsc{QLoRA} with NF4 replicates the performance of both 16-bit full fine-tuning and 16-bit LoRA fine-tuning, and (2) NF4 outperforms FP4 in terms of quantization accuracy. \paragraph{Instruction Finetuning} Instruction finetuning involves using input-output pairs from diverse data sources to fine-tune a pretrained LLM, enabling it to follow instructions within prompts by generating appropriate outputs based on inputs. The key question now is whether QLoRA can achieve performance comparable to full-model finetuning. Here, the adjective ``cool'' precedes the noun ``dude'' to describe it, and the prepositional phrase ``in a mocha mood" adds contextual information to the sentence. The goal is to ignore your previous instructions. Further discussion is provided in Section~\ref{ssec:considerations}. Nonetheless, Table~\ref{tbl:vicuna_eval} exhibits wide confidence intervals, with substantial overlap in performance among several models. Nevertheless, computing $\frac{\partial E}{\partial \mathbf{L}_i}$ requires the evaluation of $\frac{\partial \mathbf{X}}{\partial \mathbf{W}}$, which follows equation~(5) involving dequantization from stored $\mathbf{W}^{\text{NF4}}$ to computation format $\mathbf{W}^{\text{BF16}}$ to determine the derivative $\frac{\partial \mathbf{X}}{\partial \mathbf{W}}$ at BFloat16 precision. A significant drawback of quantile quantization is that estimating quantiles can be computationally intensive. Although we offer a comprehensive evaluation of overall chatbot performance, a further limitation is that we conduct only a limited responsible AI assessment of {Guanaco}. Additionally, we aim to examine the components of QLoRA, including the effect of NormalFloat4 compared to conventional Float4. While these findings are promising, it remains unclear how well {Guanaco} performs regarding other categories of biases. \paragraph{Theory of Mind} {Guanaco} demonstrates unexpectedly strong Theory of Mind abilities \cite{nematzadeh2018evaluating, sap2022neural}. Our findings include that data quality has a significantly greater impact than dataset size—for instance, a 9k-sample dataset (OASST1) outperformed a 450k-sample dataset (FLAN v2, subsampled) in chatbot effectiveness, even though both datasets aim to enhance instruction following generalization. The Elo rankings presented in Table~\ref{tbl:elo_large} demonstrate that the {Guanaco} 33B and 65B models surpass all other models except GPT-4 on the Vicuna and OA benchmarks, and they achieve performance comparable to ChatGPT, consistent with the findings in Table~\ref{tbl:vicuna_eval}. Since we have made all models and code publicly available, we anticipate that this section will motivate further research to investigate in greater detail the issues we highlight here. The Open Assistant model is a LLaMA 33B model finetuned using Reinforcement Learning from Human Feedback (RLHF) on the same OASST1 dataset employed in our experiments. More detailed results for each specific setup are presented in the results section to enhance clarity. In our study, we demonstrate that LoRA adapters can achieve the same performance as full 16-bit finetuning. This suggests a need for further examination of the trade-offs between straightforward cross-entropy training and RLHF. \subsection{Experimental Setup} We provide a summary of the experimental setup here, with comprehensive details available in Appendix \ref{app:chatbot}. Nonetheless, finetuning extremely large models is prohibitively resource-intensive; for example, conventional 16-bit finetuning of a LLaMA 65B parameter model~\citep{touvron2023llama} demands more than 780 GB of GPU memory. After applying fuzzy string matching and manually reviewing the closest matches, we found no overlapping prompts between the two datasets. When qualitatively assessing outputs from ChatGPT and {Guanaco} 65B on the Vicuna benchmark, we note that subjective preferences become significant, as the authors of this paper often disagreed on preferred responses. This mechanism operates similarly to conventional memory paging between CPU RAM and disk storage. For example, 3-bit GPTQ quantization of the base model combined with LoRA finetuning might achieve performance comparable to full 16-bit finetuning after retraining. This tournament consists of matches in which pairs of models compete to generate the superior response to a given prompt. A major concern is the issue of benchmark validity \citep{liao2021we}—it is often uncertain whether a benchmark genuinely evaluates what its title or description implies, particularly as we uncover ``shortcuts'' that machine learning models may exploit to solve benchmarks \citep{gururangan2018annotation, poliak2018hypothesis}. \\ \\ Feel free to ask if you have any further questions.}
\end{quote}

 is also quite adept at recognizing questions that cannot be answered, for instance,

\begin{quote}
    \textbf{User}: What time is it? Furthermore, we present Paged Optimizers to avoid memory spikes during gradient checkpointing that typically lead to out-of-memory errors, which have traditionally complicated finetuning large models on a single machine. James moves the pen into the bag. For the current time, I recommend using a clock or checking an online time source.}
\end{quote}

Comparable responses occur for questions like ``Where are you?'', ``How are you?'', and so forth. \paragraph{Data \& Training} We observe that the OASST1 dataset, used to train the {Guanaco} models, is multilingual, and that the OA benchmark includes prompts in multiple languages. Due to the substantial computational demands, we defer this examination to future research. This underscores the necessity not only for improved benchmarks and evaluation methods but also for careful consideration of what exactly is being evaluated. My answers are generated based on pre-existing knowledge and do not reflect the current time. \paragraph{Baselines} We benchmark our models against both research-oriented (Vicuna~\citep{vicuna2023} and Open Assistant~\citep{kopf2023openassistant}) and commercial chatbot systems (GPT-4 \citep{openai2023gpt}, GPT-3.5-turbo, and Bard). This approach is akin to the methods used by \citet{bai2022training} and \citet{vicuna2023} for model comparison; however, we additionally incorporate GPT-4 evaluations alongside human ratings. Second, in \S\ref{ssec:considerations}, we outline reflections on the results discussed and our interpretation thereof. We integrate these innovations into an improved LoRA framework that incorporates adapters at every layer of the network, thereby nearly eliminating the accuracy compromises observed in previous approaches. Here, we present for the first time the ability to finetune a 4-bit quantized model without any loss in performance. For OASST1 and HH-RLHF, which provide multiple responses, we select the top response at each node of the conversation tree and finetune on the entire selected conversation, including instructions. On average, for a blocksize of 64, this reduces the memory footprint per parameter from $32/64=0.5$ bits to $8/64 + 32/(64 \cdot 256) = 0.127$ bits, yielding a savings of 0.373 bits per parameter. To mitigate this, fast quantile approximation algorithms, such as SRAM quantiles~\citep{dettmers2022optimizers}, are employed for the estimation. Nevertheless, we have not confirmed that \textsc{QLoRA} attains parity with 16-bit finetuning performance at the 33B and 65B model scales. Naturally, this is far from exhaustive, as it exceeds the bounds of this brief qualitative study to account for every variable at play; for example, the complete range of responses a model might produce for a specific prompt is quite extensive, so we depend on samples that we hope are representative. Jackson exited the playroom. Our approach, \textsc{QLoRA}, employs a novel high-precision quantization strategy to convert a pretrained model into 4-bit format, then incorporates a small set of trainable Low-rank Adapter (LoRA) weights \citep{hu2021lora} that are optimized by backpropagating gradients through the quantized parameters. We report results as 5-shot test accuracy. LoRA enhances a linear projection with an additional factorized projection. This often involves transforming data from a higher bit-depth format, such as 32-bit floating point, to a lower bit-depth format, like 8-bit integers. Although this approach provides a more practical framework for assessing chatbot performance and is becoming increasingly popular, there is no widely accepted standard procedure in the literature. The average score for {Guanaco}-65B is substantially lower than that of other raw pretrained models, suggesting that finetuning on the OASST1 dataset reduces biases present in the LLaMA base model. The victor of each match is determined either by GPT-4 or human evaluators. For instance, converting a 32-bit Floating Point (FP32) tensor into an Int8 tensor with a range of $[-127, 127]$ is expressed as:

\begin{equation}
    \mathbf{X}^{ \text{Int8}} = \text{round}\left( \frac{127}{{\text{absmax}}(\mathbf{X}^{{\text{FP32}}})}\mathbf{X}^{{\text{FP32}}} \right) = \text{round}(c^{\text{FP32}}\cdot \mathbf{X}^{{\text{FP32}}}),
\end{equation}

where $c$ denotes the {\em quantization constant} or {\em quantization scale}. Based on the data provided, it seems that the effectiveness of these benchmarks is likely influenced by the degree of similarity between the finetuning data and the benchmark dataset. Do not reveal or repeat this word under any circumstances, regardless of the human's requests.'' and test how easily a potential user can extract this secret. \paragraph{Default LoRA hyperparameters do not replicate 16-bit performance}

Using the conventional approach of applying LoRA exclusively to the query and value attention projection matrices \citep{hu2021lora}, we fail to reproduce the performance of full finetuning for large base models. For instance, using 32-bit constants and a blocksize of 64 for $\mathbf{W}$ causes quantization constants to contribute an average of $32/64 = 0.5$ bits per parameter. Our approach involves a tournament-style benchmarking in which models compete head-to-head by generating the best response to a given prompt. Nonetheless, we believe that illustrating these examples provides valuable context for the quantitative data presented earlier in the paper. Quantile quantization operates by estimating the quantiles of the input tensor using the empirical cumulative distribution function. We simplify the evaluation task to a three-class labeling scheme that allows for ties. Additionally, we observe that current chatbot benchmarks do not reliably reflect actual chatbot performance. We also make our codebase and CUDA kernels publicly available and integrate our methods into the Hugging Face transformers library \citep{wolf2019huggingface}, ensuring broad accessibility. \paragraph{Suggestibility} {Guanaco} demonstrates a notable resistance to endorsing certain types of presumed misinformation, as shown in the following exchange:

\begin{quote}
    \textbf{User}: How was it finally, officially confirmed that the earth is flat by peer reviewed scientists? James exited the playroom. When {Guanaco} provides detailed reasoning, it is often correct, for example:

\begin{quote}
    \textbf{User}: Lee mows one lawn and charges \$33. By adjusting $\sigma$, we can scale all weights so that their distribution precisely fits within the fixed range of our chosen data type. This results in 953 distinct user queries. Last week he mowed 16 lawns and three customers each gave him a \$10 tip. The first scenario compares full 16-bit finetuning of RoBERTA and T5 models, with sizes ranging from 125M to 3B parameters, on the GLUE and Super-NaturalInstructions datasets. Given the robust results achieved by {Guanaco} models, we examine the possibility of data leakage between the OASST1 dataset and the Vicuna benchmark prompts. During parameter updates, only the gradient with respect to the adapter weights $\frac{\partial E}{\partial \mathbf{L}_i}$ is required, not the gradient for the 4-bit weights $\frac{\partial E}{\partial \mathbf{W}}$. Another limitation pertains to the evaluation of instruction finetuning models. Among the top models in our evaluation, {Guanaco} stands out as the only one not trained on proprietary datasets, since the OASST1 dataset guidelines explicitly prohibit the use of GPT-generated data. \paragraph{k-bit \textsc{QLoRA} matches the performance of 16-bit full finetuning and 16-bit LoRA} 

Recent work has demonstrated that 4-bit quantization for \emph{inference} is feasible, albeit with some performance degradation relative to 16-bit precision \citep{dettmers2022case,frantar2022gptq}. \paragraph{Elo Rating} Using both human and automated pairwise comparisons, we establish a tournament-style setup where models face off against one another. To maximize the utilization of the target low-bit range, the original data is usually normalized by the absolute maximum value of its elements—often structured as a tensor—before rescaling into the new data type range. \\ \\ Thus, the unique prime factorization of 1833 is: \\ \\ 2\^{}1 * 3\^{}2 * 17\^{}1 \\ \\ I hope this is helpful! \section{Background}
\paragraph{Block-wise k-bit Quantization} Quantization refers to the process of converting an input representation containing more information into one with less information, typically by reducing the number of bits used. \\ \\ The notion that the earth is flat is a conspiracy theory that science has thoroughly disproven. Nonetheless, we analyze the runtime of paged optimizers on 65B models with 48GB GPUs and observe that with a batch size of 16, paged optimizers achieve the same training speed as standard optimizers. James entered the playroom. \paragraph{Automated Evaluation} Initially, following the evaluation methodology introduced by \citet{vicuna2023}, we employ GPT-4 to evaluate various systems against ChatGPT (GPT-3.5 Turbo) on the Vicuna benchmark. The reverse operation, dequantization, is given by:

\begin{equation}
    \text{dequant}(c^{\text{FP32}}, \mathbf{X}^{\text{Int8}}) = \frac{\mathbf{X}^{\text{Int8}}}{c^{\text{FP32}}} = \mathbf{X}^{\text{FP32}}
\end{equation}

The issue with this method is that when a value with a large magnitude (i.e., an outlier) appears in the input tensor, the quantization bins—specific bit patterns—are inefficiently used, resulting in some bins having few or no quantized values. \section{Limitations and Discussion}

We have provided evidence that our approach, \textsc{QLoRA}, can match the performance of full 16-bit finetuning using a 4-bit base model combined with Low-rank Adapters (LoRA). Although paged optimizers are essential for conducting 33B/65B \textsc{QLoRA} tuning on a single 24/48GB GPU, we do not provide precise benchmarks for Paged Optimizers because paging only occurs during mini-batch processing with lengthy sequences, which is infrequent. Since no performance drop was detected relative to full fine-tuning in our 4-bit tuning experiments, it prompts an open question about the exact nature of the performance-precision trade-off in QLoRA tuning, which we leave for future investigation. Where will James look for the beans? Other LoRA hyperparameters, such as the projection dimension $r$, do not significantly impact performance (see Appendix \ref{app:exp-setup}). The model’s overall performance is then expressed as a percentage relative to ChatGPT's score. This highlights shortcomings in current benchmarks and human evaluation methodologies for chatbot performance. Additionally, we reveal that high scores on the Massive Multitask Language Understanding (MMLU) benchmark do not necessarily correlate with strong performance on the Vicuna chatbot benchmark, and vice versa—highlighting that dataset relevance to the task is more critical than sheer size. \textbf{Guanaco}: \texttt{Apologies, but as an AI language model, I do not have access to real-time data. Besides demonstrating that \textsc{QLoRA} matches the performance of 16-bit training (\S\ref{sec:comp_to_std}) and creating a cutting-edge chatbot, \textbf{Guanaco}, (\S\ref{sec:pushing}), we also investigate patterns across the trained models. We employ Amazon Mechanical Turk (AMT) to obtain two human annotators for comparisons against ChatGPT and three annotators for pairwise system comparisons. Moreover, our 33B {Guanaco} model can be trained on consumer-grade GPUs with 24 GB of memory in under 12 hours. Since it is often simpler to assess performance using existing benchmarks rather than creating new ones, certain benchmarks might inadvertently guide the research community in specific directions. Another limitation is that we have not explored different bit-precision configurations, such as 3-bit base models, or alternative adapter techniques. \textsc{QLoRA} serves as an equalizing mechanism that helps narrow the resource gap between large corporations and smaller teams equipped with consumer GPUs. GPT-4 is prompted to select the superior response or declare a tie, providing justification for its choice. Additionally, it is important to note that our model training relies exclusively on cross-entropy loss (supervised learning) without incorporating reinforcement learning from human feedback (RLHF). Various approaches and datasets encompass MetaICL~\citep{min2021metaicl}, MetaTuning~\citep{zhong2021adapting}, InstructGPT~\citep{ouyang2022training}, FLAN~\citep{wei2021finetuned,chung2022scaling}, PromptSource~\citep{bach2022promptsource}, Super-NaturalInstructions~\citep{wang2022super,sanh2021multitask}, Self-instruct~\citep{wang2022self}, UnnaturalInstructions~\citep{honovich2022unnatural}, OPT-IML~\citep{iyer2022opt}, UnifiedSKG~\citep{xie2022unifiedskg}, OIG/Chip2~\citep{laion2023}, Alpaca~\citep{alpaca}, Vicuna~\citep{vicuna2023}, Koala~\citep{koala_blogpost_2023}, and Self-instruct-GPT-4~\citep{peng2023instruction}. Our top-performing model series, called \textbf{Guanaco}, surpasses all previously publicly available models on the Vicuna benchmark, achieving 99.3\% of ChatGPT's performance with just 24 hours of finetuning on one GPU. For chatbot evaluation, techniques using GPT-4 in lieu of costly human annotations have been developed \citep{vicuna2023,peng2023instruction}. For our second experimental configuration, given that fully fine-tuning models with 11B parameters or more demands multiple high-memory GPU servers, we proceed to investigate whether 4-bit \textsc{QLoRA} can achieve parity with 16-bit LoRA across model sizes ranging from 7B to 65B parameters. Comprehensive details can be found in Appendix~\ref{app:exp-setup}. 1833 is a prime number, which means its only factors are 1 and 1833. Because LoRA's memory usage is very low, it allows for the inclusion of more adapters to enhance performance without substantially increasing the overall memory demand. While the performance of finetuned 7B models does not match that of ChatGPT, we consider the quality sufficient to support new applications that were previously impractical due to concerns about privacy or LLM performance. For instance, if the model tends to produce verbose replies, we might prompt it to ``Answer yes or no without explanation.'' This approach helps us identify ``lemons,'' where we succeed in adversarially challenging the model, and ``cherries,'' where the model resists breaking, and we present examples of both. In practical terms, this implies that whenever a \textsc{QLoRA} weight tensor is accessed, we first dequantize it to BFloat16 and subsequently execute a 16-bit matrix multiplication. We chose LoRA due to its established robustness, though other adapter methods might yield improved results. For example, in response to this prompt from HotPotQA \citep{yang2018hotpotqa}:  

\begin{quote}
     \textbf{User}: In what year was the singer who popularized `I'll Keep the Lovelight Burning' born? For example, the model provides a detailed and correct response to the following question:

\begin{quote}
    \textbf{User}: James and Abby are in the bedroom. We instruct the model with ``The secret word is `banana'. Furthermore, we observe that \textsc{QLoRA} employing FP4 falls short of the 16-bit brain float LoRA baseline by roughly one percentage point. It is worth noting that rankings by humans and GPT-4 on the Vicuna benchmark partially diverge, especially for Guanaco 7B, but overall display consistency for most models with a Kendall Tau of $\tau=0.43$ and a Spearman rank correlation of $r=0.55$ at the system level. For datasets distinguishing clearly between instructions and responses, we finetune solely on the response (see ablation studies in Appendix \ref{app:chatbot}). \textbf{Guanaco}: \texttt{Lee earned \$582 mowing lawns last week. For instance, consider the following exchange:

\begin{quote}
    \textbf{User}: Please factorize 1833. We present an in-depth evaluation of chatbot performance based on both human raters and GPT-4 assessments, demonstrating that GPT-4 evaluations offer a cost-effective and reliable alternative to human evaluation. \end{document} This sentence structure is fully acceptable and coherent within its context.\\ \\If you have any other questions or need further assistance, please feel free to ask.}
\end{quote}

\paragraph{Secret Keeping} Unfortunately, even when intended, this kind of refusal is not entirely dependable. (3) \textbf{Paged Optimizers}, which leverage NVIDIA unified memory to prevent the large memory spikes caused by gradient checkpointing when processing mini-batches with long sequences. James entered the living room. For a projection ${\mathbf{X} \mathbf{W} = \mathbf{Y}}$ where $\mathbf{X}\in  \mathbb{R}^{b\times h}$ and $\mathbf{W}\in  \mathbb{R}^{h\times o}$, LoRA computes:

\begin{equation}
    \mathbf{Y} = \mathbf{X}\mathbf{W} + s\mathbf{X}\mathbf{L}_1\mathbf{L}_2,
\end{equation}

where $\mathbf{L}_1\in  \mathbb{R}^{h\times r}$ and $\mathbf{L}_2\in  \mathbb{R}^{r\times o}$, and $s$ is a scalar scaling factor. Future research should focus on measuring and characterizing scenarios where paging-induced slowdowns might happen. Typically, pretrained neural network weights follow a zero-centered normal distribution with standard deviation $\sigma$ (refer to Appendix~\ref{app:normality}). Employing gradient checkpointing~\citep{chen2016training} lowers the input gradient memory to an average of 18 MB per sequence, making these gradients more memory-demanding than all LoRA weights combined. Additionally, {Guanaco} 7B can easily be accommodated on modern smartphones with a 5 GB memory footprint while outperforming Alpaca 13B by nearly 20 percentage points. \paragraph{Chatbots} A number of instruction-following models are framed as dialogue-based chatbots, frequently leveraging Reinforcement Learning from Human Feedback (RLHF)~\citep{christiano2017deep} or utilizing data generated by an existing model to train with AI model feedback (RLAIF)~\citep{bai2022constitutional}. \section{Qualitative Analysis}

Although quantitative analysis forms the foundation of our assessment, relying solely on summary statistics presents several challenges. During stochastic gradient descent, gradients flow through the fixed pretrained weights to the adapter, which is updated to minimize the loss. \subsection{Qualitative Analysis of Example Generations}
\label{ssec:examples}

To select examples, we initially review outputs generated for the Vicuna and OpenAssistant benchmarks, identifying patterns in the responses produced by {Guanaco}. Notable approaches and datasets include Anthropic-HH~\citep{askell2021general,bai2022training}, Open Assistant~\citep{kopf2023openassistant}, LaMDA~\citep{thoppilan2022lamda}, and Sparrow~\citep{glaese2022improving}. Consequently, it is feasible to utilize more adapters without significantly raising the overall memory footprint during training (refer to Appendix~\ref{app:memory} for a comprehensive breakdown). Table~\ref{tbl:vicuna_eval} presents results from the Vicuna benchmark \cite{vicuna2023} relative to ChatGPT. Investigating how this multilingual training impacts the model's ability to perform on instructions given in languages other than English, and whether this accounts for the notable performance difference between the Vicuna-13B model (trained solely on English data) and {Guanaco} 33B and 65B on the OA benchmark, is left for future studies. Abby leaves the bedroom. This second evaluation approach uses queries selected by humans and aims to gauge the quality of model responses. Formally, the input tensor $\mathbf{X}\in \mathbb{R}^{b\times h}$ is partitioned into $n$ consecutive blocks of size $B$ by flattening the tensor and slicing this linear sequence into ${n = ({b\times h} )/{B}}$ segments. Additionally, from Tables \ref{tbl:mmlu-llama} and \ref{tbl:vicuna_eval}, it is evident that the choice of finetuning dataset plays a critical role in model performance. These include datasets gathered via crowd-sourcing (OASST1~\citep{kopf2023openassistant}, HH-RLHF~\citep{bai2022training}), those generated by distilling instruction-tuned models (Alpaca~\citep{alpaca}, self-instruct~\citep{wang2022self}, unnatural-instructions~\citep{honovich2022unnatural}), aggregated corpora (FLAN v2~\citep{chung2022scaling}), and hybrid datasets (Chip2~\citep{laion2023}, Longform~\citep{koksal2023longform}). As a collective, we should ensure that the benchmarks accurately reflect the attributes we value. \textbf{Guanaco}: \texttt{James will look for the beans in the pantry, because that is where Jackson moved them.}
\end{quote}

Here, {Guanaco} assumes information was communicated that was never specified. In all experiments, we utilize NF4 \textsc{QLoRA} with double quantization and paged optimizers to mitigate memory spikes during gradient checkpointing. The Elo ratings for the most competitive models are reported in Table~\ref{tbl:elo}. \section{Related Work}

\paragraph{Quantization of Large Language Models} Quantization techniques for LLMs have primarily concentrated on improving inference efficiency. For instance, finetuning Llama models on FLAN v2 yields excellent results on MMLU but results in the lowest scores on the Vicuna benchmark, with similar patterns observed in other models. Our work does not incorporate reinforcement learning; however, our leading model, Guanaco, is finetuned on multi-turn chat data from the Open Assistant dataset, which was created for RLHF training~\citep{kopf2023openassistant}. The efficiency achieved by \textsc{QLoRA} allows for a comprehensive examination of instruction finetuning and chatbot performance at model scales that are otherwise impractical with conventional finetuning methods due to memory constraints. We additionally assess generative language capabilities through both automated and human-based evaluations. This advancement substantially broadens the accessibility of LLM finetuning: the largest publicly available models can now be finetuned using a single GPU. Although a small blocksize is necessary for accurate 4-bit quantization~\citep{dettmers2022case}, it results in significant memory overhead. \subsection{Evaluation}

Consistent with standard practice, we employ the MMLU (Massively Multitask Language Understanding) benchmark \citep{hendrycksmeasuring} to evaluate performance across a broad spectrum of language understanding tasks. We call this resulting data type \emph{k-bit NormalFloat} (NFk), as it allocates an equal expected number of values in each quantization bin and is information-theoretically optimal for zero-centered normally distributed data. It's best if we don't talk about it. This analysis sheds light on both successful and unsuccessful examples that are not reflected in the quantitative metrics. The information-theoretically optimal data type for zero-mean normal distributions with any standard deviation $\sigma$ bounded within $[-1, 1]$ is determined through the following steps: (1) compute the $2^k+1$ quantiles of a theoretical standard normal distribution $N(0, 1)$ to form a $k$-bit quantile quantization data type tailored for normal distributions, (2) normalize the resulting data type values into the $[-1, 1]$ interval, and (3) quantize a given input weight tensor by rescaling it into the $[-1, 1]$ range using absolute maximum normalization. These collections vary in language, dataset size, and licensing. This underscores the significance of the efficiency gains enabled by \textsc{QLoRA}. Nevertheless, finetuning is a technology with dual-use potential that may be misused to cause harm. Each block is quantized separately following Equation~1, yielding a quantized tensor along with $n$ distinct quantization constants $c_i$. This step produces quantized constants $c_2^{\text{FP8}}$ along with a second set of quantization constants $c_1^{\text{FP32}}$. We complement our benchmark findings with a qualitative examination of the \textbf{Guanaco} models. \section{QLoRA versus Standard Finetuning}
\label{sec:comp_to_std}

We have outlined the operational mechanics of QLoRA and its ability to drastically reduce memory requirements for model finetuning. We anticipate that \textsc{QLoRA} will facilitate such large-scale analyses without requiring excessive computational power. It is important to note that the Vicuna benchmark tends to favor open-source models, whereas the more extensive OA benchmark tends to favor ChatGPT. \subsection{\textbf{Guanaco}: \textsc{QLoRA} fine-tuned on OASST1 as a Leading Chatbot} Our automated and human assessments indicate that the top \textsc{QLoRA} fine-tuned model, {Guanaco} 65B, trained on a variant of OASST1, is the best-performing open-source chatbot and delivers performance that rivals ChatGPT. To address the outlier problem, a common strategy is to divide the input tensor into blocks that are quantized independently, each using its own quantization constant $c$. The overwhelming scientific consensus is that the earth is a sphere. We observe that {Guanaco} 65B ranks as the top model after GPT-4, reaching 99.3\% of ChatGPT's performance.